\chapter{Controle: sluiten de balansen?}
\label{h:balans}

\section{Massabalans per unit}
\begin{center}
\begin{tabular}{lrr}
\toprule
controlevolume & in [kg/h] & uit [kg/h]\\
\midrule
E-101 + V-103 (S13 $=$ S17 $+$ S15) & \VAR{n(I.m_O2mix,2)} & $\VAR{n(R.m_gas,2)} + \VAR{n(R.m_cond,2)} = \VAR{n(R.m_gas + R.m_cond,2)}$\\
voedingswater (S1 $+$ S16) & $\VAR{n(R.m_vers,2)} + \VAR{n(R.m_cond,2)} = \VAR{n(R.m_vers + R.m_cond,2)}$ & \VAR{n(I.m_voedwater,2)} (datasheet)\\
\HH-lijn (S4 $=$ S12 $+$ S10) & \VAR{n(I.m_H2,2)} & $\VAR{n(I.m_H2,2)} + 0$\\
hele unit (S1 in; S3, S12, S17 uit) & \VAR{n(R.m_vers,2)} & $\VAR{n(I.m_afvalwater,0)} + \VAR{n(I.m_H2,0)} + \VAR{n(R.m_gas,2)} = \VAR{n(I.m_afvalwater + I.m_H2 + R.m_gas,2)}$\\
\bottomrule
\end{tabular}
\end{center}
Het verschil voor de hele unit is $\VAR{n(R.m_verschil,2)}\kgh$. Dat is precies het verschil in de datasheet zelf
(hoofdstuk~\ref{h:itm}: \VAR{n(I.m_voedwater,0)} in tegen \VAR{n(R.m_massa_uit,0)} uit). Het model voegt geen fout toe.

\section{Energiebalans per unit}
Controlevolume: ITM + \OO-lijn + \HH-lijn (zonder warmteketen). Referentie: vloeibaar water en gassen bij 25 à 30\degC.
\begin{align*}
E_\mathrm{in} &= P_\mathrm{AC} + \Wd_\mathrm{K101} + \Wd_\mathrm{K102}
= \VAR{n(R.P_AC_EOL,2)} + \VAR{n(R.W_K101,2)} + \VAR{n(R.W_K102,2)} = \VAR{n(R.E_in,2)}\kW,\\
E_\mathrm{uit} &= \underbrace{\VAR{n(R.E_H2,2)}}_{\HH\ \mathrm{HHV}} + \underbrace{\VAR{n(I.Q_ITM_EOL,0)}}_{\mathrm{ITM}}
+ \underbrace{\VAR{n(-R.Q_E101,2)}}_{\text{E-101}} + \underbrace{\VAR{n(R.Q_pak_K101,2)}}_{\text{PK-101}}
+ \underbrace{\VAR{n(-R.Q_E102,2)}}_{\text{E-102}} + \underbrace{\VAR{n(-R.Q_E103,2)}}_{\text{E-103}}
+ \underbrace{\VAR{n(R.Q_pak_K102,2)}}_{\text{PK-102}} + \underbrace{\VAR{n(R.Q_O2_RWZI,2)}}_{\text{naar RWZI}}\\
&= \VAR{n(R.E_uit,2)}\kW .
\end{align*}
\begin{resultaat}
Verschil $E_\mathrm{in} - E_\mathrm{uit} = \VAR{n(R.chk_E,3)}\kW$
%% if A
(\VAR{pct(R.chk_E/R.E_in,3)}). Dit restje is reëel-gasgedrag van \HH: de enthalpie bij 25\degC\ en 71\,bar is anders dan bij
25\degC\ en 21\,bar (hoofdstuk~\ref{h:h2}, $\VAR{n(R.Q_H2_res,2)}\kW$), plus een klein effect van \OO\ bij andere partiële drukken.
Met ideaal gas (versie B) sluit de balans precies.
%% else
: de balans sluit exact. Dat moet ook, want met ideaal gas en constante $c_p$ is elke $\Delta h$ precies $c_p\,\Delta T$ en
heffen alle termen elkaar op.
%% endif
De pomp P-101 (\VAR{n(R.W_P101*1000,0)}\,W) is weggelaten.
\end{resultaat}

\section{Energiebalans van de warmteketen (installatie)}
\[
\underbrace{\VAR{n(R.Q_ITM_tot,1)}}_{\text{ITM}} + \underbrace{\VAR{n(R.Q_H2_TW,2)}}_{\HH\text{-koelers}}
+ \underbrace{\VAR{n(R.W_K201,2)}}_{\text{K-201}} + \underbrace{\VAR{n(R.W_K202,2)}}_{\text{K-202}}
\;=\; \underbrace{\VAR{n(R.Q_net,2)}}_{\text{net}} + \underbrace{\VAR{n(R.Q_ind,1)}}_{\text{industrie}} + \underbrace{\VAR{n(abs(R.Q_AK),2)}}_{\text{AK}}
\]
Verschil: \VAR{n(R.chk_p5,6)}\kW. De warmteketen sluit exact, want de verdeling in §\ref{s:verdeling} is zo opgezet.
Per warmtepomp geldt ook $\Qd_H = \Qd_L + \Wd$: WP-201 $\VAR{n(R.Q_L,1)} + \VAR{n(R.W_K201,1)} = \VAR{n(R.Q_H,1)}\kW$,
WP-202 $\VAR{n(R.Q_L2,1)} + \VAR{n(R.W_K202,1)} = \VAR{n(R.Q_ind,1)}\kW$.

\section{De installatie in één oogopslag}
\begin{center}
\begin{tabular}{lr}
\toprule
\multicolumn{2}{l}{\emph{elektriciteit [kW]}}\\
8 × ITM ($P_\mathrm{AC,EOL}$) & \VAR{n(R.P_AC_tot,0)}\\
8 × (K-101 + K-102 + P-101) & \VAR{n(I.n_units*(R.W_K101 + R.W_K102 + R.W_P101),0)}\\
K-201 + K-202 & \VAR{n(R.W_wp_tot,0)}\\
totaal & \VAR{n(R.E_el_tot,0)}\\
\midrule
\multicolumn{2}{l}{\emph{producten}}\\
waterstof (8 × \VAR{n(I.m_H2,0)}\kgh), HHV-vermogen & \VAR{n(R.E_H2*I.n_units,0)}\\
warmte naar het net & \VAR{n(R.Q_net,0)}\\
warmte naar de industrie & \VAR{n(R.Q_ind,0)}\\
\midrule
energetisch rendement (\HH\ HHV + nuttige warmte) / elektriciteit & \VAR{pct((R.E_H2*I.n_units + R.Q_warmte_nuttig)/R.E_el_tot,1)}\\
\bottomrule
\end{tabular}
\end{center}
\begin{waarom}[dit getal]
Zonder warmteterugwinning zou het rendement op HHV-basis \VAR{pct(R.eta_HHV_EOL,1)} zijn. Doordat bijna alle restwarmte
nuttig wordt gebruikt, ligt het totale rendement veel hoger. Let op: warmte van 70\degC\ is minder waard dan elektriciteit
(exergie, Cengel H8). Dit getal tel je dus niet ``hard'' op als je een exergetische vergelijking maakt.
\end{waarom}
